\par\SourceLeaf{492}%
Im bevorstehenden Sommer wird er zuerst an der Nordspitze Jütlands die astronomischen Beobachtungen anstellen. Danach wird er in Lauenburg mit neuen Instrumenten die vorjährigen Messungen wiederholen und die Azimute bestimmen sowie bei Hamburg eine einige Meilen lange Grundlinie messen. Meine Anwesenheit wird nur für einen Teil der Arbeiten in Lauenburg und für die Basismessung gewünscht. Für die Arbeiten an der Nordspitze Jütlands wird sie nicht gewünscht. Die dänische Regierung scheint mir mit ihrem Wunsch nach meiner Anwesenheit bei diesen heiklen Arbeiten davon ausgegangen zu sein, dass eine gemeinsame Beratung bei dieser allein zum Wohl der Wissenschaft unternommenen Arbeit förderlich sein werde. Durch bloßen Briefwechsel ließe sich das der Natur der Sache nach nur unvollkommen oder gar nicht erreichen. Zweitens soll die ganze Gradmessung künftig in einem eigenen Werk in vollem Umfang der gelehrten Welt bekannt gemacht werden. Die Anwesenheit eines Zeugen bei einigen der wichtigsten Arbeiten wird daher dazu beitragen, die Authentizität der Gradmessung zu stärken. So lud Frankreich im Jahr 1798 zur Basismessung bei Paris alle Staaten, die sich damals mit Frankreich nicht im Krieg befanden, ein, geeignete Astronomen zu entsenden. Dänemark, die Schweiz, Holland, Spanien und andere Staaten taten dies. Diese beiden Erwägungen beziehen sich zunächst nur auf die dänischen Gradmessungen für sich genommen. Unter der Voraussetzung, dass früher oder später eine weitere Fortsetzung durch Hannover angeordnet werden könnte, scheint es aber noch besonders wichtig, dass ein Astronom von unserer Seite das Vertrauen gebührend würdigen kann, das diejenigen Arbeiten besonders verdienen, an die eine solche Fortsetzung unmittelbar anschließen müsste.

Nach diesen Überlegungen bin ich sehr gern bereit, der für mich schmeichelhaften Aufforderung zu folgen. Ich werde den dänischen Astronomen bitten, mich möglichst früh darüber zu unterrichten, wann diese Arbeiten vorgenommen werden, damit ich meine übrigen Angelegenheiten danach einrichten kann. Bei den Vorlesungen wird das wohl keine Schwierigkeiten bereiten, da die Arbeiten erst im Spätsommer beginnen werden. Es wird mir zwar leidtun, die in ihrer Art einzigartige große ringförmige Sonnenfinsternis am 7. September aller Wahrscheinlichkeit nach hier nicht mehr beobachten zu können. Ich werde mich aber damit trösten, dass sie hier doch von Prof. \textsc{Harding} beobachtet werden kann und dass sie auch in Holstein ringförmig erscheinen wird. Dort wird es leicht möglich sein, für die Beobachtung dieses bemerkenswerten Phänomens gute Hilfsmittel zu beschaffen, wenn sie auch den hiesigen nicht ganz gleichkommen. — — —

Göttingen, den 27. Februar 1820.
\begin{flushright}C. F. \textsc{Gauss}.\end{flushright}

\par\bigskip\Needspace{8\baselineskip}
\pdfbookmark[0]{P. M. über die hannoversche Gradmessung.}{pm-hannoversche-gradmessung}
\begin{center}\large P. M. ÜBER DIE HANNOVERSCHE GRADMESSUNG.\end{center}
\par
— — — In meinem früheren Bericht habe ich bereits dargelegt, dass mehrere Hauptinstrumente für diese Arbeit unentbehrlich sind. Solange noch nicht einmal feststand, woher sie zu bekommen wären, konnten die eigentlichen Arbeiten natürlich noch nicht beginnen, abgesehen von einigen astronomischen Vorarbeiten an der hiesigen Sternwarte, mit denen ich schon im vergangenen Sommer angefangen habe. Inzwischen habe ich diese wichtigsten Geräte, namentlich einen größeren, möglichst vollkommenen Theodoliten und ein sogenanntes Universalinstrument, bei dem bedeutendsten heute lebenden Instrumentenmacher bestellt, bei von \textsc{Reichenbach} in München. Er hat die Bestellung angenommen und versprochen, die Instrumente Anfang Mai beziehungsweise im Juli dieses Jahres zu liefern.
